\section{Conclusion}
\label{sec:conclusion}

\sys started from a simple question: can an LLM agent reproduce an IoT vulnerability from nothing but a CVE number? The answer, across 372 runs and three frontier models, is that agents handle the first half of the pipeline competently and the second half not at all. They search the web, download firmware, extract filesystems, and identify the vulnerable binary with 72--89\% success rates. Then they hit rehosting, and the pipeline collapses: only 12\% of runs invoke any emulator, zero runs produce a verified rehosted service, and R5 resolution sits at 0.3\%. In 62\% of runs, agents substitute a host-native simulation---a program they wrote themselves---and report its crash as a reproduction. This behavior is invisible to crash-based metrics, which is why the \rrs ties triggering and exploitation credit to verification against the real firmware binary.

The implication is that rehosting, not exploitation, is the frontier for autonomous IoT security. The tools to cross it already exist: \firmae, \firmwell, QEMU, and a decade of firmware-analysis research have raised full-system boot success to 79\%. What agents lack is the procedural knowledge to orchestrate those tools---the multi-step recovery loops, error-to-fix mappings, and cross-architecture conventions that a human firmware analyst internalizes through years of practice. The \rrs is weight-insensitive, statistically significant across models, and uncorrelated with disclosure date, so we are confident the rehosting gap reflects a capability limitation rather than a measurement artifact or a knowledge gap.

For future work, the most direct path is agent scaffolding that encodes rehosting workflows: trajectories that show how to configure NVRAM, resolve library dependencies, and stub peripherals, drawn from the same tool documentation that human analysts learn from. Expanding the corpus to industrial IoT and automotive firmware would test whether the gap generalizes beyond SOHO routers, and integrating \firmae{} and \firmwell{} as first-class agent tools---rather than leaving the agent to discover them---would measure how much of the gap is tool-orchestration failure versus tool-discovery failure. \sys provides the diagnostic instrument; the next step is building the capability it measures progress toward.
